\chapter{MATERIAIS E MÉTODOS}
\label{capitulo2}

\section{MATERIAIS}

Descreva os materiais utilizados na pesquisa.

\section{MÉTODOS}

Descreva a metodologia geral da pesquisa.

\subsection{Metodologia DNER (1998)}

Descreva a metodologia DNER utilizada.

\begin{table}[htbp]
  \centering
  \caption{Índices para soma ponderada}
  \label{tab:indices_ponderada}
  \small
  \begin{tabular}{lc}
    \hline
    \textbf{Tipo de Acidente} & \textbf{Índice} \\
    \hline
    Com vítima fatal & 13 \\
    Com ferido & 3 \\
    Sem vítima & 1 \\
    \hline
  \end{tabular}
  \\[0.2cm]
  {\small Fonte: DNER (1998).}
\end{table}

\subsection{Metodologia Vadeby e Forsman (2018)}

Descreva a metodologia Vadeby e Forsman.

\subsubsection{Subtítulo do 3º nível}

O 3º nível pode ser itemizado por letras (a, b, c...) ou por números romanos (i, ii, iii...). Uma vez escolhida a forma de identificação dos itens, este deve ser mantido para todo o texto. Apenas o primeiro e o segundo níveis são numerados no sumário.

\subsection{Teste t de Student}

Descreva o teste estatístico utilizado.
